\chapter{La máquina completa}
% versión completa: chapters/10_synthesis.tex

\pencilfig{10}{\pencilart{10}}{La máquina completa: el bus de datos, los frenos dentro de él y cuatro canales de radio de control encima.}

Tomamos los tipos de neuronas de uno en uno y nos preguntamos qué le agrega cada uno a la máquina. Es hora de juntar las piezas.

Resultó una máquina de tres capas. Abajo, la enorme red telefónica de neuronas \nt{GLUT}: por ella viajan los datos. Al lado, las neuronas inhibidoras \nt{GABA}: controlan el flujo, deciden qué pasa, cuándo y con qué volumen. Y por encima de todo, unas pocas emisoras de radio, cada una con su perilla. \nt{DOPA} dice qué cambios de conexiones conservar. \nt{SERO} mantiene el nivel general y, según la hipótesis, el horizonte de paciencia. \nt{NORA} elige entre buscar y decidir. \nt{ACET}, entre escribir y leer. En realidad, cada emisora es un pequeño haz de canales y no un solo cable, pero aun así los canales son un puñado para ochenta y seis mil millones de neuronas.

\section*{Una sola historia}

Sigamos un suceso a través de toda la máquina. Un animal sabe desde hace tiempo que si aprieta la palanca A recibe jugo. Un día el mundo cambia: A ya no funciona, y el jugo lo da la palanca B, que el animal casi no usaba.

Tras A no llegó el jugo: la dopamina responde con un silencio, y las conexiones que eligieron A se debilitan. Los silencios se repiten, los errores se vuelven más grandes que de costumbre: la noradrenalina, según la hipótesis, avisa que el mundo cambió. La perilla de contraste baja, la elección se vuelve más blanda, y el ruido empuja más a menudo a probar B. La acetilcolina pasa la red a modo escritura. Probar B trae un jugo inesperado: la dopamina estalla y las conexiones que eligieron B se refuerzan. Tras unos cuantos intentos, la expectativa alcanza al mundo nuevo, las sorpresas se acaban y las perillas vuelven a su sitio.

Fíjese: ninguna perilla sabe lo que hacen las demás. Cada una responde a sus propias señales, y el orden de las acciones se arma solo, como en un circuito asíncrono donde cada bloque actúa cuando sus entradas están listas. Que cada perilla se comporte así por separado está, en lo esencial, medido. Que trabajen juntas de forma coordinada es por ahora una hipótesis.

\section*{En qué no se parece a su portátil}

La portátil tiene un reloj común; el cerebro, no. Los elementos de la portátil son fiables; una sinapsis pierde más de la mitad de las señales. La memoria de la portátil está separada del procesador; en el cerebro está en las mismas conexiones. La portátil aprende (si aprende) por retropropagación del error; el cerebro, con reglas locales y una única señal de difusión “mejor de lo esperado”. Y lo principal: la portátil gasta decenas de vatios en un solo procesador, y el cerebro, veinte vatios en todo.

\section*{Lo que no sabemos}

No sabemos cuánto aprovecha la corteza el tiempo exacto de los clics. No sabemos cómo ajusta miles de millones de conexiones en circuitos de muchos pisos si tiene un solo maestro para todas. No entendemos del todo qué transmiten las emisoras. No sabemos cómo coordinan su trabajo partes lejanas del cerebro sin un reloj común. Y por ahora nadie sabe construir una máquina que haga lo mismo con veinte vatios.

A partir de aquí, el libro sale de este núcleo: hacia las entradas y salidas de la máquina, su depuración, el sueño y las máquinas de neuronas vivas sobre un chip.

\fullversion{10}
